%======================================================================
\section{Evolution of the Dark Sector}
\label{sec:evolution}
%======================================================================

The cosmological evolution of the secluded dark sector can be naturally divided into two stages, each amenable to a distinct formulation.
The first stage tracks the coupled evolution of the DM particle $X$, the mediator $Y$, and the dark-sector temperature as the dark sector departs from internal chemical equilibrium. The key processes during this phase are the $2\to2$ annihilation $X\bar{X}\to YY$, which sets the DM relic abundance, and the $3\to2$ cannibal reactions, which maintain the mediator in chemical equilibrium until they too freeze out. Throughout this work, the portal coupling $\varepsilon$ is assumed to be sufficiently small that the hidden sector remains thermally decoupled from the Standard Model bath, evolving independently at its own temperature $T_h = \xi\,T_{\rm SM}$. In this regime, the mediator decay width generically satisfies $\Gamma_Y \ll H$ during the entire freeze-out process, so that it can be safely neglected and the dynamics is governed solely by the dark-sector couplings. The evolution during this stage is then captured by a set of coupled Boltzmann equations, which we present in Sec.~\ref{sec:boltzmann}. We first write these equations in their standard, physically transparent form and then show how they can be recast in terms of chemical-potential variables and the hidden to SM temperature ratio $\xi$, a formulation that is better suited for numerical integration.

The second stage begins once the Boltzmann system has converged and the comoving abundances of $X$ and $Y$ are effectively frozen. The relevant dynamics is now the decay of the long-lived mediator into SM particles. As $Y$ depletes, it injects entropy into the visible sector and dilutes the dark-matter yield; it is through this mechanism that the portal coupling $\varepsilon$ ultimately controls the predicted relic abundance. This late-time background evolution is described in Sec.~\ref{sec:background}.

\subsection{Boltzmann Equations}
\label{sec:boltzmann}

The starting point is the standard set of Boltzmann equations for the
number densities $n_X$ and $n_Y$ and the total dark-sector energy
density $\rho_h = \rho_X + \rho_Y$.  In a Friedmann--Robertson--Walker background these read:
\begin{align}
    &\dot{n}_X + 3H n_X = -\sigmav_{X}\left(n_X^2 - (n_X^{\text{eq}})^2\,\frac{n_Y^2}{(n_Y^{\text{eq}})^2}\right) \label{eq:boltz_nX} \\
    &\dot{n}_Y + 3H n_Y = \sigmav_{X}\left(n_X^2 - (n_X^{\text{eq}})^2\,\frac{n_Y^2}{(n_Y^{\text{eq}})^2}\right) - \Gamma_{\rm can}\,(n_Y - n_Y^{\text{eq}}) \label{eq:boltz_nY} \\
    &\dot{\rho}_h + 3H(\rho_h + P_h) = 0 \label{eq:boltz_rho}
\end{align}
where a dot corresponds to a derivative with respect to cosmic time, $H$ is the total Hubble rate including the dark-sector contribution to the energy density, Eq.~\eqref{eq:Htot}, $\sigmav_{X}$ is the thermally averaged annihilation cross section, and $\Gamma_{\rm can}$ collects all $3\to2$ number-changing processes, namely:
\begin{equation}
  \Gamma_{\rm can} = \langle\sigma v^2\rangle_{YYY\to YY}\,n_Y^2\,+\langle\sigma v^2\rangle_{YYX\to YX}\,n_Y\,n_X+\,
\langle\sigma v^2\rangle_{YXX\to XX}\,n_X^2\,.
\label{eq:Gamma_can_std}
\end{equation}


The first equation governs the DM abundance, with the collision term vanishing when $n_X/n_Y = n_X^{\rm eq}/n_Y^{\rm eq}$ to enforce detailed balance. The second equation tracks the mediator, which is sourced by annihilation and depleted by cannibal reactions. The third equation expresses energy conservation within the dark sector and determines the evolution of the dark-sector temperature $T_h$. Under the Maxwell--Boltzmann statistics  adopted throughout this work, the energy density and pressure of each species are fully specified by its number density and temperature through $\rho_i = B_1^i\,n_i$ $\rho_i + P_i =B_2^i\, n_i$, Eqs.~\eqref{eq:MB_rho_P} and~\eqref{eq:B1}, so that Eq.~\eqref{eq:boltz_rho} can be rewritten as an evolution equation for the hidden-sector temperature $T_h$ given $n_X$ and $n_Y$.

\subsubsection{Reformulation in terms of chemical potentials and $\xi$}
\label{sec:chempot_formulation}

Rather than evolving the number densities directly, it is advantageous to reformulate the system in terms of the dimensionless chemical potentials $\bar{\mu}_i \equiv \mu_i/\TD$, defined through
\begin{equation}
    Y_i = Y_i^{\text{eq}}\,e^{\bar{\mu}_i}\,, \qquad i \in \{X, Y\}\,,
    \label{eq:chempot_def}
\end{equation}
where $Y_i = n_i/s$ is the comoving yield, $s$ is the SM entropy density and $Y_i^{\text{eq}} = \neqMB_i/s$ its equilibrium value. This parametrisation has two key benefits. First, in the early universe when both species are in chemical equilibrium, $\bar{\mu}_i = 0$ exactly, avoiding the need to track the cancellation between $Y_i$ and $Y_i^{\text{eq}}$ --- a severe numerical challenge when both are large and nearly equal. Second, during freeze-out the dimensionless chemical potential grows linearly while the yield departs exponentially, so $\bar{\mu}_i$ is a much more slowly varying quantity and hence better suited as an ordinary differential equation (ODE) variable.

We take as independent variable $x \equiv m_X/T_{\rm SM}$ and promote
the dark-to-SM temperature ratio $\ln\xi$ to a dynamical degree of
freedom, obtaining a coupled system of three ODEs:
\begin{align}
    \frac{d\bar{\mu}_X}{dx} &= A + B\,\frac{d\ln\xi}{dx} \label{eq:dmuX} \\
    \frac{d\bar{\mu}_Y}{dx} &= C + D\,\frac{d\ln\xi}{dx} \label{eq:dmuY} \\
    \frac{d\ln\xi}{dx} &= \frac{E + F\,A + G\,C}{1 - F\,B - G\,D} \label{eq:dlnxi_full}
\end{align}
where the six coefficient functions $A$--$G$ depend on $(x,\,\xi,\,\bar\mu_X,\,
\bar\mu_Y)$ and are constructed from the thermally averaged cross
sections, the equilibrium thermodynamic quantities ($B_1$, $B_2$,
$\lambda$), and the SM entropy-conservation factor $\tilde g$ as defined in Secs.~\ref{sec:cosmology} and~\ref{sec:models}. The
explicit forms for these coefficients are given in Eqs.~\eqref{eq:Acoeff}--\eqref{eq:Gcoeff}
below.
 
\paragraph{Dark-matter coefficients.}
%
\begin{align}
  A &= -\frac{s\,\tilde g}{H_{\rm tot}\,x}\,\langle\sigma v\rangle_{X} Y_X^{\rm eq}e^{\bar\mu_X}
       \left[1- e^{2(\bar\mu_Y-\bar\mu_X)}
       \right]
       + \frac{\lambda_X - 3\tilde g}{x}\,,
  \label{eq:Acoeff}
  \\
  B &= -\lambda_X\,.
  \label{eq:Bcoeff}
\end{align}
%
The first term in $A$ is the collision integral for $X\bar X \leftrightarrow YY$, written in a form that makes the approach to chemical equilibrium ($\bar\mu_X = \bar\mu_Y$) manifest. The second term encodes the redshifting of the equilibrium yield, with $\lambda_X \equiv d\ln n_X^{\rm eq}/d\ln T_h$, Eq.~\eqref{eq:lambda}, capturing the temperature sensitivity of the equilibrium density and $\tilde g$, Eq.~\eqref{eq:gtilde}, which accounts for the running of the SM entropy degrees of freedom.
 
\paragraph{Mediator coefficients.}
%
\begin{align}
  C &= \frac{s\,\tilde g}{H_{\rm tot}\,x}\,\left\{
       \langle\sigma v\rangle_{X}\, \frac{(Y^{\rm eq}_X)^2}{Y^{\rm eq}_Y}e^{2\bar\mu_X-\bar\mu_Y}\,\left[1-e^{2(\bar\mu_Y-\bar\mu_X)}\right]
       - \Gamma_{\rm can}\bigl(1 - e^{-\bar\mu_Y}\bigr)\right\}
       + \frac{\lambda_Y - 3\tilde g}{x}\,,
  \label{eq:Ccoeff}
  \\
  D &= -\lambda_Y\,,
  \label{eq:Dcoeff}
\end{align}
The first contribution in $C$ is the annihilation source, equal and opposite to the corresponding term in $A$. The second involves the total cannibal reaction rate $\Gamma_{\rm can}$ defined in Eq.~\eqref{eq:Gamma_can_std}, multiplied by the factor $(1 - e^{-\bar\mu_Y})$ which vanishes identically when $Y$ is in chemical equilibrium ($\bar\mu_Y = 0$) and acts as a restoring force that drives $\bar\mu_Y$ back toward zero whenever cannibal reactions are fast. The last term plays the same role as in $A$, accounting for the redshifting of $Y_Y^{\rm eq}$.
 
\paragraph{Temperature-equation coefficients.}
These arise from projecting the energy-conservation equation onto
$d\ln\xi/dx$:
%
\begin{align}
  E &= \frac{1}{x}
       \left(1 - 3\tilde g\,
             \frac{n_X\,B_2^X + n_Y\,B_2^Y}{\tilde D}
       \right),
  \label{eq:Ecoeff}
  \\
  F &= -\frac{B_1^X\,n_X}{\tilde D}\,,
  \qquad
  G = -\frac{B_1^Y\,n_Y}{\tilde D}\,,
  \label{eq:Gcoeff}
\end{align}
%
with the energy denominator
\begin{equation}
  \tilde D
  = T_h\!\left(
      n_X\,\frac{dB_1^X}{dT_h}
    + n_Y\,\frac{dB_1^Y}{dT_h}
    \right)
  + B_1^X\,n_X\,\lambda_X
  + B_1^Y\,n_Y\,\lambda_Y\,.
  \label{eq:Dtilde}
\end{equation}
%
Here, the coefficient $E$ captures the adiabatic cooling of the hidden sector due to Hubble expansion, while $F$ and $G$ encode the feedback of changes in $n_X$ and $n_Y$ on the hidden-sector temperature. All three are controlled by the denominator $\tilde D$, which measures the total heat capacity of the dark plasma. The thermodynamic quantities $B_1$, $B_2$, and $\lambda$ used here are defined
in Sec.~\ref{sec:bessel}.


\section{Details on the numerical implementation and alternative solvers}
\label{app:numerics}

This appendix collects additional technical details on the numerical integration of the Boltzmann system presented in Sec.~\ref{sec:boltzmann}. We first describe three alternative solver strategies that are retained in the codebase for testing and cross-validation purposes but are not used in production. We then summarize the numerical settings common to all solvers, including the integration grid, tolerances, and regularization procedures.

\subsection{Alternative solvers}
\label{app:legacy}

Three additional solvers are implemented alongside the primary three-phase solver described in Sec.~\ref{sec:solver}. They differ in how aggressively the phase structure is exploited and in the choice of dynamical variables during the intermediate epoch.

The two-phase solver (\texttt{solve\_boltzmann\_chempot\_2phase}) skips the intermediate Phase~1.5 entirely, transitioning directly from full equilibrium (Phase~1) to the complete three-variable system (Phase~2) when $\Gamma_{\rm ann}/H$ drops below $10^7$. This introduces unnecessary stiffness near the annihilation freeze-out transition, where $Y$ is still effectively in equilibrium but $X$ is departing, precisely the difficulty that the intermediate phase of the primary solver is designed to resolve.

The hybrid solver (\texttt{solve\_boltzmann\_hybrid}) does include an intermediate phase, but formulates it in terms of the physical yields $Y_X$ and $Y_Y$ rather than the dimensionless chemical potentials:
\begin{align}
\frac{dY_X}{dx} &= -\frac{s\,\tilde g}{H_{\rm tot}\,x}\,\sigmav_X\left(Y_X^2 - (Y_X^{\rm eq})^2\,\frac{Y_Y^2}{(Y_Y^{\rm eq})^2}\right), \label{eq:dYX_hybrid} \\
\frac{dY_Y}{dx} &= -\frac{dY_X}{dx} - \frac{\tilde g}{H_{\rm tot}\,x}\,\Gamma_{\rm can}\,(Y_Y - Y_Y^{\rm eq})\,,
\label{eq:dYY_hybrid}
\end{align}
where $\Gamma_{\rm can}$ is the cannibal rate defined in Eq.~\eqref{eq:Gamma_can_std}.

It transitions to the chemical-potential variables of Phase~2 when either $\Gamma_{\rm can}/H$ falls below a threshold or the fractional departure $|Y_X - Y_X^{\rm eq}|/Y_X^{\rm eq}$ exceeds a configurable value. While conceptually straightforward, tracking $Y_Y$ directly can suffer from numerical loss of precision when it spans many orders of magnitude during the cannibal epoch.

The $Y$-equilibrium solver (\texttt{solve\_boltzmann\_Yeq}) is the simplest of the four, assuming that the mediator always tracks its equilibrium abundance throughout the entire evolution. It solves only two variables, $(\ln\xi,\,\bar\mu_X)$, using equations identical to Phase~1.5 of the primary solver. By construction, this solver cannot capture the cannibal freeze-out of $Y$ or its subsequent departure from equilibrium, which affects both the late-time temperature evolution and the mediator abundance at the onset of the background evolution stage.

\subsection{Numerical settings}
\label{app:num_details}

All solvers use \texttt{scipy.integrate.solve\_ivp} with the implicit Runge--Kutta method Radau (L-stable, suitable for stiff problems). The maximum step size is set to $1$ in $x$, with relative tolerance $10^{-8}$ and absolute tolerance $10^{-15}$ during Phase~2 (relaxed to $10^{-12}$ during Phase~1). Rather than integrating from $x_{\min}$ to $x_{\max}$ in a single call, the evolution is broken into segments over a predefined grid of breakpoints,
\begin{equation}
    x \in \{10^{-2},\, 10^{-1},\, 0.99\} \cup \mathrm{linspace}(1,\,10,\,20) \cup \{10.5,\, 11.5,\, \ldots,\, 99.5\} \cup \{200,\, 500,\, \ldots,\, x_{\rm max}\}\,,
\end{equation}
with $500$ evaluation points per segment. This segmented approach improves solver robustness: the integrator restarts fresh at each breakpoint, which helps it adapt its internal step size to the rapidly changing stiffness of the system. Phase-switching conditions are nevertheless checked continuously within each segment.

Several clipping and regularisation procedures are applied to prevent numerical overflow and underflow. The temperature ratio is restricted to $\ln\xi \in [-20,\,20]$ and the dimensionles chemical potentials to $|\bar\mu_i| \leq 10/x$ per Hubble time. Exponentials of yields are clipped to the range $[-700,\,700]$, and densities that fall below $10^{-300}$ are replaced with $\pm 10^{-300}$, preserving the original sign.

In addition to the default convergence criterion of the primary solver described in Sec.~\ref{sec:solver}, an alternative mode is available in which the solver instead monitors the rate of change of the DM yield and terminates when $|d\ln Y_X/d\ln x|$ falls below the configurable tolerance (the parameter \texttt{convergence\_threshold} in the code). This mode can be selected via the \texttt{convergence\_mode} option.
